\section{Methodology in detail}

\begin{figure}[htbp]
\centering
\begin{tikzpicture}[node distance=0.6cm and 0.9cm,
  box/.style={draw=#1,fill=#1!8,rounded corners=3pt,align=center,inner sep=5pt,font=\small,text width=3.0cm},
  box/.default=blue, arr/.style={-{Stealth[length=2mm]},thick}]
  \node[box=gray] (data) {Historical data\\(training period)};
  \node[box,right=of data] (fit) {Fit AR(1) dynamics\\ + innovation covariance};
  \node[box,right=of fit] (graph) {Correlation graph\\ $\Acorr$};
  \node[box=black,below=of graph] (sink) {Kalman sink\\ estimate $\hat x_t$, covariance $P_t$};
  \node[box=green!50!black,left=of sink] (guard) {Freshness guard\\ forced set $\mathcal F_t$};
  \node[box,left=of guard] (gnn) {GNN scheduler\\ picks $C-|\mathcal F_t|$ more};
  \node[box=orange,below=of guard] (net) {Sensors transmit\\ $y_t = x_t[S_t]+n_t$};
  \draw[arr] (data)--(fit); \draw[arr] (fit)--(graph); \draw[arr] (graph)--(sink);
  \draw[arr] (sink)--(guard); \draw[arr] (guard)--(gnn);
  \draw[arr] (gnn) |- (net); \draw[arr] (net) -| (sink);
  \draw[arr,dashed] (graph.north) -- ++(0,0.4) -| node[pos=0.25,above,font=\scriptsize]{graph and features} (gnn.north);
\end{tikzpicture}
\caption{The complete pipeline. The model and graph are fitted once on training data; the loop at the bottom runs
every slot.}
\end{figure}

\subsection{System model}

\begin{itemize}
  \item $M$ sensors with fixed positions. Time is divided into slots $t=1,2,\dots$ (5 minutes on real data).
  \item In each slot the sink chooses a set $S_t$ of at most $C$ sensors (\textbf{capacity constraint}).
  \item Optionally, sensors closer than an \textbf{interference radius} cannot transmit in the same slot (used only in
        the synthetic interference study; real-data experiments have no interference).
  \item Polled sensors send their reading with a small measurement noise: $y_t = x_t[S_t] + n_t$,
        $n_t\sim\mathcal N(0,rI)$, $r=10^{-3}$.
  \item Every sensor's age must satisfy $a_i(t)\le F$ (\textbf{freshness constraint}).
  \item Slot timeline: (1) the field moves one step, (2) the sink predicts, (3) the freshness guard computes forced
        sensors, (4) the scheduler completes the poll set, (5) the polled sensors transmit, (6) the sink updates,
        ages are updated, and the error is measured.
\end{itemize}

\subsection{Field model and the synthetic simulator}

Readings are modelled as a stationary, first-order autoregressive process with correlated innovations:
\[
  x_t = \operatorname{diag}(\rho)\,x_{t-1} + w_t,\qquad w_t\sim\mathcal N(0,Q).
\]
$\rho_i\in[0,1)$ says how much of sensor $i$'s value carries over to the next slot (temporal correlation); the
off-diagonal entries of $Q$ say how the sensors' surprises move together (spatial correlation).

\textbf{Synthetic simulator.} To test on thousands of controlled layouts, \texttt{CorrelatedField} generates fields
with a known ground truth:
\begin{enumerate}
  \item Place $M$ sensors in the unit square: \emph{uniform} (spread evenly) or \emph{clustered} (around 2--7 random
        centres with spread 0.07).
  \item Spatial correlation by an RBF kernel: $\Sigma_{ij}=\exp\!\big(-\lVert p_i-p_j\rVert^2/(2\ell^2)\big)$, where
        $\ell$ (length scale, default 0.3) controls how far correlation reaches.
  \item Per-sensor speed $\rho_i = \bar\rho \pm$ spread (default $\bar\rho=0.9$).
  \item Innovations $Q = D\Sigma D$ with $D=\operatorname{diag}(\sqrt{1-\rho_i^2})$, so every sensor has unit variance.
        The stationary covariance is $P_{0,ij}=Q_{ij}/(1-\rho_i\rho_j)$.
\end{enumerate}
The \textbf{layout seed} (positions, $\Sigma$, $\rho$) and the \textbf{noise seed} are separate. This fixes the
mid-semester bug: the graph is always estimated on the same layout the scheduler runs on (a unit test checks the
estimated and true correlations agree with $r>0.95$). Every policy is evaluated on the \emph{identical} trajectory
(the field is rewound between policies), so differences come only from the scheduling.

\subsection{Fitting the model from data}

The sink is never told the true model; it fits one from data (\texttt{fit\_model}). Given a block of training
readings $x_1,\dots,x_T$ (standardised):
\begin{align*}
  \hat\rho_i &= \frac{\sum_t x_{t-1,i}\,x_{t,i}}{\sum_t x_{t-1,i}^2}\quad(\text{clipped to }[0,0.9995]),\\
  e_t &= x_t - \hat\rho \odot x_{t-1}\quad(\text{innovations}),\qquad
  \hat Q = \operatorname{cov}(e_t)+10^{-6}I,\qquad \hat P_0=\operatorname{cov}(x_t)+10^{-6}I.
\end{align*}
\begin{whybox}[Why one-step least squares?]
It is the maximum-likelihood fit of an AR(1) model and is cheap even for 325 sensors and months of data. We also tried
fitting the autocorrelation at longer lags; on the real datasets it made the sink 4--17$\times$ over-confident (its
reported uncertainty was far below its actual error), so one-step fitting was kept.
\end{whybox}
For synthetic layouts the sink fits on a \textbf{pilot phase} of 1000 slots in which it hears every sensor. For real
data it fits on the training period only.

\subsection{The correlation graph}

From $\hat Q$ we compute the innovation correlation $\hat r_{ij}=\hat Q_{ij}/\sqrt{\hat Q_{ii}\hat Q_{jj}}$ and keep
strong edges:
\[
  (\Acorr)_{ij} = |\hat r_{ij}|\cdot\mathbf 1\big[|\hat r_{ij}|\ge 0.3\big],\qquad (\Acorr)_{ii}=0.
\]
\begin{whybox}[Why innovations, not raw readings?]
Raw readings of almost all sensors are correlated simply because they share trends (every room warms during the day,
every freeway slows at rush hour). That would connect everything to everything. Innovations remove the predictable
part, so an edge means ``when one sensor is surprised, the other is surprised in the same way'': exactly the
information that lets the sink infer one from the other. The threshold 0.3 keeps meaningful edges and drops noise.
\end{whybox}

\subsection{The Kalman sink}

The sink holds a mean $m$ and covariance $P$:
\begin{align*}
  \textbf{Predict:}\quad & m^- = \rho\odot m,\qquad P^- = D_\rho P D_\rho + Q,\qquad D_\rho=\operatorname{diag}(\rho)\\
  \textbf{Update:}\quad & K = P^-_{:,S}\,\big(P^-_{SS}+rI\big)^{-1},\quad m = m^- + K\,(y - m^-_S),\quad
  P = P^- - K\,P^-_{S,:}
\end{align*}
Here $S$ is the set of polled sensors. Polling sensor $j$ reduces $P_{ii}$ for every $i$ with $P_{ij}\neq0$: the
correlated neighbours learn too (a unit test checks this). $P$ does not depend on the actual readings, only on
\emph{who} was polled, so the scheduler can plan with it.

\textbf{Variance gain.} If sensor $i$ alone were polled next, the total uncertainty would drop by
\[
  g_i = \frac{\sum_k P_{ki}^2}{P_{ii}+r}.
\]
This single number drives the Kalman-greedy baseline and the first term of our decoder.

\begin{whybox}[Why a Kalman filter, not a neural forecaster?]
It is optimal for the linear-Gaussian model, needs no training, gives a calibrated uncertainty $P$ (a unit test
checks that the predicted variance matches the actual error within 8\%), and updates in milliseconds. A learned
forecaster would not give us $P$, which is what makes correlation-aware scheduling possible.
\end{whybox}

\subsection{The freshness guard}

Before any scheduler acts:
\begin{enumerate}
  \item Find all sensors with $a_i(t)\ge F$.
  \item Sort them oldest first and take up to $C$ of them that do not conflict with each other: this is the forced
        set $\mathcal F_t$.
  \item The scheduler only fills the remaining $C-|\mathcal F_t|$ places.
  \item The environment checks (asserts) in every slot that capacity, conflicts and freshness hold.
\end{enumerate}
We use $F=\lfloor 3M/C\rfloor$: three times the round-robin period. Intel: $F=\lfloor 3\cdot45/6\rfloor=22$ slots
(1.8 hours). PEMS-BAY: $F=\lfloor3\cdot325/40\rfloor=24$ slots (2 hours). Synthetic headline: $F=15$.

\begin{whybox}[Why a hard guard instead of a penalty in the reward?]
A penalty can always be traded away: if ignoring a sensor reduces error enough, the policy will ignore it. A guard
makes the guarantee unconditional and identical for every scheduler, which also makes comparisons fair.
Section 8 shows what happens without it: Kalman-greedy never polls 18 of 50 sensors.
\end{whybox}

\begin{whybox}[Why $F = 3M/C$?]
Round robin needs $M/C$ slots to visit everyone. A bound equal to $M/C$ would force pure round robin and leave no
freedom; a very large bound would allow near-starvation. Three times the period gives the scheduler real freedom
while still keeping every sensor within a couple of hours on real data.
\end{whybox}

\subsection{Objective and MDP formulation}

\textbf{Objective.}
\[
  \min_\pi\ \mathbb E\Big[\frac{1}{TM}\sum_{t=1}^T\lVert x_t-\hat x_t\rVert^2\Big]
  \quad\text{s.t.}\quad |S_t|\le C,\ \text{no conflicts in } S_t,\ a_i(t)\le F\ \ \forall i,t.
\]

\textbf{As an MDP:}
\begin{itemize}
  \item \textbf{State}: ages $a(t)$, which sensors were polled last slot, the sink's covariance $P_t$, the learned
        $\hat\rho$ and graph $\Acorr$, the conflict matrix, the forced set.
  \item \textbf{Action}: an ordered list of $C-|\mathcal F_t|$ sensors to poll.
  \item \textbf{Transition}: the field evolves, the sink predicts and updates, ages change.
  \item \textbf{Reward}: $r_t = -(1-\gamma)\log(\text{error}_t + 0.005)$ (details in 5.9).
\end{itemize}

\subsection{The GNN scheduler}

\subsubsection{Node features (13 per sensor)}
\begin{center}
\small
\begin{tabularx}{\textwidth}{@{}rlX@{}}
\toprule
\# & Feature & Why it is there \\
\midrule
1 & $\min(a_i/(M/C),6)/6$ & Age relative to the round-robin period (comparable across network sizes) \\
2 & polled last slot (0/1) & Recently polled sensors usually need no repeat \\
3 & $P_{ii}$ & The sink's uncertainty about sensor $i$ \\
4 & $\log(P_{ii}+10^{-3})/7$ & Same on a log scale (uncertainty spans orders of magnitude) \\
5 & $g_i/M$ & How much total uncertainty polling $i$ removes \\
6 & $\log(g_i/M+10^{-4})/9$ & Same on a log scale \\
7 & forced (0/1) & Already committed by the guard \\
8 & slack $=\mathrm{clip}((F-a_i)/F,0,1)$ & How soon the guard will force this sensor \\
9 & $C/M$ & How scarce the channel is \\
10 & $\hat\rho_i$ & How persistent the sensor is \\
11 & $\log(1-\hat\rho_i^2+10^{-5})/12$ & How fast it drifts away from its last value \\
12 & degree in $\Acorr$ / $(M-1)$ & How connected (redundant) it is \\
13 & guard on (0/1) & Whether a freshness bound applies \\
\bottomrule
\end{tabularx}
\end{center}
All features are scaled to roughly $[0,1]$ so one network works across network sizes.

\subsubsection{Two graphs}
\begin{itemize}
  \item $\Acorr$: the static correlation graph (who is correlated in general).
  \item $\Apost$ with $(\Apost)_{ij}=|P_{ij}|/\sqrt{P_{ii}P_{jj}}$: the sink's \emph{current} posterior correlation,
        recomputed every slot (who is redundant \emph{right now}, given what was just polled).
\end{itemize}
Both are row-normalised ($\tilde A = A/\text{row sum}$) before use.

\subsubsection{Encoder}
\begin{align*}
  h^{(0)}_i &= \mathrm{MLP}_{\text{in}}(\text{features}_i)\in\mathbb R^{64}\\
  h^{(\ell+1)} &= \mathrm{LN}\Big(h^{(\ell)} + \mathrm{MLP}_\ell\big[\,h^{(\ell)},\ \tilde\Acorr h^{(\ell)},\ \tilde\Apost h^{(\ell)},\ \bar h^{(\ell)}\,\big]\Big),\quad \ell=0,1,2
\end{align*}
$\bar h$ is the mean over all sensors (global context). Each MLP is $4d\to2d\to d$ with GELU. Three layers let
information travel three hops. The residual connection and LayerNorm keep training stable.

\subsubsection{Residual pointer decoder}
Sensors are picked one at a time. At each step the score (logit) of sensor $i$ is
\[
  \text{logit}_i = \underbrace{\alpha\log g_i}_{\text{Kalman-greedy}} + \underbrace{f_\theta(h_i)}_{\text{learned correction}}
  + \underbrace{\textstyle\sum_{j\in\text{picked}} G_{ij}}_{\text{learned interaction}},\qquad G_{ij}=\frac{q_i^\top k_j}{\sqrt{32}}.
\]
\begin{enumerate}
  \item Mask sensors already picked, forced, or in conflict with a pick.
  \item Softmax over the remaining scores; \emph{sample} a sensor during training, take the \emph{argmax} at test time.
  \item Condition $P$ on that pick (Kalman update), recompute every $g_i$, add $G_{\cdot,\text{pick}}$.
  \item Repeat until the slot is full.
\end{enumerate}

\textbf{Residual initialisation.} The output layers of $f_\theta$ and $q$ are initialised to zero, so at the start
$f_\theta=G=0$ and the decoder \emph{is} Kalman-greedy (with $\alpha=4$ the softmax is sharp). A unit test checks that
the untrained network makes exactly the same choices as Kalman-greedy for 40 slots.

\textbf{Plackett--Luce probability.} The probability of the ordered pick list is the product of the softmax
probabilities at each step, so PPO gets exact log-probabilities (also unit-tested).

Size: 137{,}282 parameters. The same weights run for any $M$.

\begin{whybox}[Why a residual decoder?]
Our first attempt trained a GNN ``student'' to imitate Kalman-greedy (behaviour cloning) and then fine-tune; it never
fully matched the teacher. Building Kalman-greedy \emph{into} the network means training starts from the strongest
model-based scheduler, and RL only has to learn \emph{when deviating from greedy pays off later}, for example polling a
sensor before the guard forces it at a bad moment, or not trusting a wrong model.
\end{whybox}

\subsection{Training with PPO}

\begin{algorithm}[htbp]
\caption{PPO training of the GNN scheduler}
\begin{algorithmic}[1]
\State Actor $\pi_\theta$ = Kalman-greedy at init; critic $V_\phi$ (2-layer GNN)
\For{iteration $=1..400$}
  \State Sample a batch of environments: synthetic (random $M\in[10,100]$, $C$, $\ell$, $\rho$, layout, $F$, interference)
  \Statex\hspace{1.3em}or, with probability 0.35, replays of a real dataset's training period
  \For{$t=1..96$ slots}
    \State $\mathcal F_t\gets$ guard; sample picks from $\pi_\theta$; step; reward $r_t=-(1-\gamma)\log(\text{err}_t+0.005)$
  \EndFor
  \State Advantages with GAE ($\gamma=0.95$, $\lambda=0.9$)
  \For{4 epochs over minibatches of 512}
    \State ratio $\eta=\pi_\theta/\pi_{\text{old}}$; loss $-\min(\eta\hat A,\ \mathrm{clip}(\eta,0.85,1.15)\hat A) - 0.002\cdot\text{entropy}$
    \State Adam step (lr $10^{-4}$ actor, $5\times10^{-4}$ critic, grad-norm clip 0.5)
  \EndFor
  \State Every 10 iterations: validate, keep the best checkpoint
\EndFor
\end{algorithmic}
\end{algorithm}

\textbf{Training scenarios.} Each synthetic batch draws: $C=\mathrm{round}(M\cdot U[0.08,0.35])$; length scale
log-uniform in $[0.05,0.45]$; $\bar\rho\in U[0.6,0.995]$; per-sensor spread with probability 0.6; clustered layout with
probability 0.4; freshness bound $3M/C$ with probability 0.7 (else off); interference with probability 0.2. The
number of parallel environments shrinks for large $M$ ($32\cdot(40/M)^2$, clipped to 8--32) to keep memory constant.

\textbf{Reward.} For synthetic environments the score is the sink's true posterior variance (computed with the true
model); for real replays it is the actual squared error against the recorded readings. The logarithm puts small
networks (large errors) and large networks (small errors) on the same scale, and the $(1-\gamma)$ factor keeps
returns near the reward scale.

\textbf{Real-data replays.} On a real dataset's training period, the sink's model is fitted on the first 70\%, the
episodes run on the next 20\% (on compact sub-networks: a random sensor and its nearest neighbours), and the last
10\% is validation. Because the model is fitted on different days than the episode, the policy experiences the
model mismatch it will face at test time.

\textbf{Validation and selection.} Validation uses four fixed synthetic scenarios (default $M=20,C=5,F=12$;
clustered; heterogeneous $\rho$; $M=60,C=8,F=22,\ell=0.15$) with layout seeds 9000+, plus real validation slices.
The checkpoint with the best validation score is kept. Test layouts use seeds 10000+ and real test periods; neither
is ever seen in training.

\textbf{Run.} 400 PPO iterations (resumable; warm-started from an earlier run), about 35 minutes on a laptop CPU.
The critic is warmed up for 15 iterations before the actor starts updating.

\subsection{Baselines and evaluation protocol}

\begin{center}
\small
\begin{tabularx}{\textwidth}{@{}lX@{}}
\toprule
\textbf{Scheduler} & \textbf{Rule (after the forced sensors)} \\
\midrule
Random & Pick uniformly at random. Lower reference. \\
Max-AoI (B1) & Pick the oldest sensors (ties by index, so round robin). AoI-based scheduling \cite{2,16}. \\
Corr-diverse (v4) & Pick the oldest; then multiply every other sensor's score by $(1-|\hat r|)$ with the picked one; repeat. Our mid-semester heuristic. \\
Max-uncertainty & Pick the largest $P_{ii}$, ignoring overlap between picks. \\
Kalman-greedy & Repeatedly pick the largest $g_i$ and condition $P$ on it \cite{19}. Strongest model-based baseline. \\
Oracle-greedy & Kalman-greedy using the \emph{true} covariance (synthetic only; not deployable). \\
GNN-PPO (ours) & Learned policy (greedy decode). \\
Ablation & Our model retrained with the message passing removed. \\
\bottomrule
\end{tabularx}
\end{center}

\textbf{Metrics.}
\begin{itemize}
  \item \textbf{MSE}: mean over sensors and slots of $(x-\hat x)^2$ in standardised units.
  \item \textbf{RMSE in units}: the same error converted back to $^\circ$C or mph (what the faculty can relate to).
  \item \textbf{Mean AoI} and \textbf{max AoI}; \textbf{starved} = number of sensors never polled.
  \item \textbf{Decision time} per slot (milliseconds on CPU).
\end{itemize}
The first 50 slots of each synthetic episode are burn-in and not scored.

\textbf{Statistics.} Mean $\pm$ 95\% CI ($1.96\cdot\text{sd}/\sqrt n$) over test layouts or windows. Because every
scheduler runs on the identical readings, we compare them \emph{paired}: the \textbf{win rate} is the share of
windows where the GNN had lower error, and the \textbf{Wilcoxon signed-rank test} gives the p-value.
